% Part: turing-machines
% Chapter: machines-computations
% Section: halting-states

\documentclass[../../../include/open-logic-section]{subfiles}

\begin{document}

\olfileid[hi]{tur}{mac}{hal}
\olsection{विराम-अवस्थाएँ}

\begin{explain}
यद्यपि हमने अपनी मशीनों को केवल तभी रुकने के लिए परिभाषित किया है जब
निष्पादित करने हेतु कोई निर्देश न हो, ट्यूरिंग मशीनों के सामान्य
निरूपणों में एक समर्पित \emph{विराम-अवस्था}~$h$ होती है, ऐसी कि
$h \in Q$।

विराम-अवस्था के पीछे का विचार सरल है: जब मशीन अपनी क्रिया पूरी कर चुकी
होती है (वह आगत स्वीकार करने को तैयार होती है अथवा निर्गत लिख चुकी होती
है), तब वह अवस्था~$h$ में जाती है और वहाँ रुकती है। कुछ मशीनों में दो
विराम-अवस्थाएँ होती हैं---एक आगत को स्वीकार करती है और दूसरी उसे
अस्वीकार करती है।
\end{explain}

\begin{ex}\emph{विराम-अवस्थाएँ}।
इस धारणा को स्पष्ट करने के लिए सम मशीन में एक परिवर्तन से आरंभ करें।
आगत सम होने पर मशीन को अवस्था~$q_0$ में रुकवाने के स्थान पर हम मशीन को
किसी विराम-अवस्था में भेजने वाला निर्देश जोड़ सकते हैं।
\[
\begin{tikzpicture}[->,>=stealth',shorten >=1pt,auto,node distance=2.8cm,
                    semithick]
  \tikzstyle{every state}=[fill=none,draw=black,text=black]

  \node[initial, state]         (A)                     {$q_0$};
  \node[state]         (B) [right of=A] {$q_1$};
  \node[state]         (C) [below of=A] {$h$};

  \path (A) edge [bend left] node {\TMtrans{\TMstroke}{\TMstroke}{R}} (B)
  	    edge node {\TMtrans{\TMblank}{\TMblank}{N}} (C)
        (B) edge [loop above] node {\TMtrans{\TMblank}{\TMblank}{R}} (B)
            edge [bend left] node {\TMtrans{\TMstroke}{\TMstroke}{R}} (A);
\end{tikzpicture}
\]

उदाहरण को और बढ़ाएँ। जब मशीन निर्धारित करती है कि आगत विषम है, तब वह
कभी नहीं रुकती। चक्रकारी निर्देश को अस्वीकार-अवस्था~$r$ में जाने वाले
निर्देश से बदलकर मशीन में एक \emph{अस्वीकार-अवस्था} सम्मिलित कर सकते
हैं।
\[
\begin{tikzpicture}[->,>=stealth',shorten >=1pt,auto,node distance=2.8cm,
                    semithick]
  \tikzstyle{every state}=[fill=none,draw=black,text=black]

  \node[initial,state]         (A)                     {$q_0$};
  \node[state]         (B) [right of=A] {$q_1$};
  \node[state]         (C) [below of=A] {$h$};
  \node[state]         (D) [below of=B] {$r$};

  \path (A) edge [bend left] node {\TMtrans{\TMstroke}{\TMstroke}{R}} (B)
  	    edge node {\TMtrans{\TMblank}{\TMblank}{N}} (C)
        (B) edge node {\TMtrans{\TMblank}{\TMblank}{N}} (D)
            edge [bend left] node {\TMtrans{\TMstroke}{\TMstroke}{R}} (A);
\end{tikzpicture}
\]
\end{ex}

\begin{explain}
इस प्रकार की स्थितियों में समर्पित विराम-अवस्था जोड़ना लाभकारी हो सकता
है, क्योंकि उससे यह स्पष्ट होता है कि मशीन किन आगतों को स्वीकार या
अस्वीकार करती है। किंतु यह ध्यान रखना महत्त्वपूर्ण है कि समर्पित
विराम-अवस्था जोड़ने से संगणन-शक्ति में कोई वृद्धि नहीं होती। इसी प्रकार,
रुकने की कम आकारिक धारणा के अपने लाभ हैं। इस अध्याय में अब तक प्रयुक्त
रुकने की परिभाषा \emph{विराम समस्या} के प्रमाण को सहज और प्रदर्शित करने
में सरल बनाती है। इसी कारण हम अपनी मूल परिभाषा जारी रखते हैं।
\end{explain}

\end{document}
